\SourceLeaf{14}
\begin{center}
{\Large\textsc{ARITHMETICAL INVESTIGATIONS.}}

\BookRule{\linewidth}

\textsc{FIRST SECTION}

{\scriptsize\textsc{ON}}

\textsc{THE CONGRUENCE OF NUMBERS IN GENERAL.}

\BookRule{0.10\linewidth}
\end{center}

\Needspace{5\baselineskip}
\begin{center}
\emph{Congruent numbers, moduli, residues and nonresidues.}
\end{center}

\Article{1}

If the number $a$ divides the difference between the numbers $b, c$, $b$ and $c$ are said to be \emph{congruent with respect to $a$}; otherwise, \emph{incongruent}: we call $a$ itself the \emph{modulus}. Each of the numbers $b, c$ is called a \emph{residue} of the other in the former case, but a \emph{nonresidue} in the latter.

These notions apply to all integers, both positive and negative\HistoricalNote{*)}{The modulus must always be taken \emph{absolutely}, that is, without any sign.}, but are not to be extended to fractions. \emph{For example,} $-9$ and $+16$ are congruent with respect to the modulus $5$; $-7$ is a residue of $+15$ with respect to the modulus $11$, but a nonresidue with respect to the modulus $3$. Also, since every number divides zero, every number is regarded as congruent to itself modulo any modulus.

\Article{2}

All the residues of a given number $a$ with respect to the modulus $m$ are contained in the formula $a+km$, where $k$ denotes an indeterminate integer. The simpler propositions that we will give later can be proved from this without difficulty; but anyone can see their truth just as easily by inspection.

\SourceLeaf{15}
From now on we will denote the congruence of numbers by the sign $\equiv$, adding the modulus in parentheses where necessary, $-16\equiv 9\ (\mathrm{mod.}\ 5)$. $-7\equiv 15\ (\mathrm{mod.}\ 11)$\HistoricalNote{*)}{We have adopted this sign because of the close analogy between equality and congruence. For the same reason, the distinguished Le Gendre, in the paper we will often cite below, kept the equality sign itself for congruence; we have hesitated to follow him, to avoid ambiguity.}.

\Article{3}

\textsc{Theorem.} \emph{Given $m$ successive integers}
\[
a,\ a+1,\ a+2\ \GaussFourDots\ a+m-1
\]
\emph{and another integer $A$, exactly one of the former is congruent to the latter modulo $m$.}

If $\frac{a-A}{m}$ is an integer, then $a\equiv A$; if it is a fraction, let the next greater integer (or, when it is negative, the next \emph{smaller}, disregarding the sign) be $=k$. Then $A+km$ will fall between $a$ and $a+m$, and so is the required number. It is clear that all the quotients $\frac{a-A}{m}$, $\frac{a+1-A}{m}$, $\frac{a+2-A}{m}$ etc. lie between $k-1$ and $k+1$; so no more than one of them can be an integer.

\Needspace{5\baselineskip}
\begin{center}\emph{Least residues.}\end{center}
\Article{4}

Every number will therefore have a residue both in the series $0, 1, 2,\ldots m-1$ and in the series $0, -1, -2,\ \GaussFourDots\ -(m-1)$, which we will call \emph{least residues}. It is clear that, unless $0$ is a residue, there are always two, one \emph{positive}, the other \emph{negative}. If they are unequal in magnitude, one will be $<\frac{m}{2}$; otherwise both will be $=\frac{m}{2}$, without regard to sign. So every number has a residue not exceeding half the modulus, which will be called the \emph{absolutely least residue}.

\emph{E.g.} $-13$, modulo $5$, has the least positive residue $2$, which is also the absolutely least residue, while $-3$ is the least negative residue; $+5$, modulo $7$, is its own least positive residue, while $-2$ is the least negative and also the absolutely least residue.

\Needspace{5\baselineskip}
\begin{center}\emph{Elementary propositions concerning congruences.}\end{center}
\Article{5}

With these notions established, let us collect the properties of congruent numbers that are immediately apparent.

\SourceLeaf{16}
\emph{Numbers congruent modulo a composite number are also congruent modulo any divisor of it.}

\emph{If several numbers are congruent to the same number modulo the same modulus, they will be congruent to one another} (modulo that same modulus).

The modulus is also understood to be the same in the statements that follow.

\emph{Congruent numbers have the same least residues; incongruent numbers have different ones.}

\Article{6}

\emph{If there are any number of integers $A, B, C$ etc. and as many others $a, b, c$ etc., congruent to them modulo any modulus,}
\[
A\equiv a,\ B\equiv b\ \text{\itshape etc.,\quad then}\quad A+B+C+\text{\itshape etc.}\equiv a+b+c+\text{\itshape etc.}
\]

\emph{If $A\equiv a, B\equiv b$, then $A-B\equiv a-b$.}

\Article{7}

\emph{If $A\equiv a$, then also $kA\equiv ka$.}

If $k$ is a positive number, this is just a particular case of the proposition in the preceding article, with $A=B=C$ etc., $a=b=c$ etc. If $k$ is negative, $-k$ will be positive, so $-kA\equiv -ka$, and then $kA\equiv ka$.

\emph{If $A\equiv a, B\equiv b$, then $AB\equiv ab$.} For $AB\equiv Ab\equiv ba$.

\Article{8}

\emph{If there are any number of integers $A, B, C$ etc. and as many others $a, b, c$ etc. congruent to them, $A\equiv a, B\equiv b$ etc., the products of each set will be congruent, $ABC\ \text{\itshape etc.}\equiv abc\ \text{\itshape etc.}$}

By the preceding article, $AB\equiv ab$, and for the same reason $ABC\equiv abc$; in the same way, any number of further factors may be added.

If all the numbers $A, B, C$ etc. are taken to be equal, and likewise the corresponding numbers $a, b, c$ etc., we have this theorem: \emph{If $A\equiv a$ and $k$ is a positive integer, then $A^{k}\equiv a^{k}$.}

\Article{9}

\emph{Let $\mathrm{X}$ be an algebraic function of the indeterminate $x$, of the form}
\[
Ax^{a}+Bx^{b}+Cx^{c}+\text{\itshape etc.}
\]
\emph{where $A, B, C$ etc. denote any integers, and $a, b, c$ etc. denote non-negative integers. Then, if values congruent modulo any modulus are assigned to the indeterminate $x$, the resulting values of the function $\mathrm{X}$ will be congruent.}

\SourceLeaf{17}
Let $f, g$ be congruent values of $x$. Then, by the preceding article, $f^{a}\equiv g^{a}$ and $Af^{a}\equiv Ag^{a}$, and in the same way $Bf^{b}\equiv Bg^{b}$ etc. Hence
\[
Af^{a}+Bf^{b}+Cf^{c}+\text{etc.}\equiv Ag^{a}+Bg^{b}+Cg^{c}+\text{etc.}\quad\text{\itshape Q. E. D.}
\]

It is also easy to see how to extend this theorem to functions of several indeterminates.

\Article{10}

So, if all consecutive integers are substituted for $x$, and the values of the function $\mathrm{X}$ are reduced to least residues, these will form a series in which the same terms recur after an interval of $m$ terms (where $m$ denotes the modulus); in other words, the series will be formed from a \emph{period} of $m$ terms repeated infinitely often. Let, for example, $\mathrm{X}=x^{3}-8x+6$ and $m=5$; then, for $x=0, 1, 2, 3$ etc., the values of $\mathrm{X}$ yield the least positive residues $1, 4, 3, 4, 3, 1, 4$ etc., where the first five, $1, 4, 3, 4, 3$, repeat indefinitely. And if the series is continued backwards, that is, negative values are assigned to $x$, the same period appears with the order of its terms reversed. This makes it clear that no terms other than those in this period can occur anywhere in the series.

\Article{11}

In this example, then, $\mathrm{X}$ can be neither $\equiv 0$ nor $\equiv 2$ (mod. 5), much less $=0$ or $=2$. It follows that the equations $x^{3}-8x+6=0$ and $x^{3}-8x+4=0$ cannot be solved by integers and consequently, as is well known, cannot be solved by rational numbers. In general it is evident that the equation $\mathrm{X}=0$, where $\mathrm{X}$ is a function of the unknown $x$, of the form
\[
x^{n}+Ax^{n-1}+Bx^{n-2}+\text{etc.}+N
\]
with $A, B, C$ etc. integers and $n$ a positive integer (a form to which, as is known, all algebraic equations can be reduced), has no rational root if the congruence $\mathrm{X}\equiv 0$ cannot be satisfied modulo some modulus. But this criterion, which has arisen naturally here, will be treated more fully in Sect. VIII. This example can certainly give at least some idea of the usefulness of these investigations.

\SourceLeaf{18}
\Needspace{5\baselineskip}
\begin{center}\emph{Some applications.}\end{center}
\Article{12}

Many things usually taught in arithmetic books rest on the theorems given in this chapter, for example the rules for testing whether a proposed number is divisible by 9, 11, or other numbers. \emph{Modulo} 9 all powers of the number 10 are congruent to one. Thus, if the proposed number has the form $a+10b+100c+\text{etc.}$, it will give the same least residue modulo 9 as $a+b+c+\text{etc.}$ So, if the individual digits of a number written in decimal notation are added without regard to the positions they occupy, this sum and the proposed number give the same least residues. So the given number is divisible by 9 if the sum is divisible by 9, and conversely. The same also holds for the divisor 3. Since, \emph{modulo} 11, $100\equiv 1$, we have in general $10^{2k}\equiv 1$. $10^{2k+1}\equiv 10\equiv -1$, and a number of the form $a+10b+100c+\text{etc.}$ gives the same least residue modulo 11 as $a-b+c$ etc.; the familiar rule follows at once. All similar rules are readily deduced from the same principle.

The preceding results also explain the rules usually recommended for checking arithmetic operations. When other numbers are to be obtained from given numbers by addition, subtraction, multiplication, or raising to powers, the given numbers are replaced by their least residues modulo an arbitrary modulus (usually 9 or 11, since in our decimal system residues for these moduli can be found so easily, as we have just shown). The numbers arising from this must be congruent to those derived from the proposed numbers; if they are not, we conclude that an error has crept into the calculation.

But these and similar matters are well known, so there is no need to spend more time on them.

\par\bigskip\bigskip\BookRule{0.29\linewidth}
